% فصل ۳ — روش کار
% همهٔ عددها از docs/generated/thesis_numbers.json می‌آیند، که src/thesis_numbers.py آن را
% از آرشیو runs/ و docs/generated/ می‌سازد. صحت با «python3 tools/farsi_lint.py --numbers» بررسی می‌شود.

\chapter{روش کار}
\label{ch:method}

این فصل تصمیم‌های روشی را توصیف می‌کند. سهم اصلی این پژوهش روش‌شناختی است، نه کارایی؛ پس
ترتیب فصل نیز همین را دنبال می‌کند: نخست پروتکل ارزیابی، بعد داده، بعد معماری، و در پایان
سازوکارهایی که تبار داده را تضمین می‌کنند. معماری آخرین چیزی است که اهمیت پیدا می‌کند، و
فصل \ref{ch:results} نشان می‌دهد چرا.

\section{پروتکل ارزیابی}

پروتکل پیش از نخستین اجرای واقعی نوشته و منجمد شد. تغییر آن نیازمند افزودن بخشی تاریخ‌دار در
انتهای پرونده است که بگوید چه چیزی و چرا عوض شد، و هر نتیجه‌ای که زیر نسخهٔ پیشین محاسبه شده
برچسب خودش را نگه می‌دارد. این قاعده وجود دارد چون تغییر خاموش پروتکل برای بهتر کردن ظاهر یک
عدد، دقیقاً همان شکستی است که این پرونده برای جلوگیری از آن نوشته شده است.

\subsection{واحد استقلال}

واحد استقلال بیمار است، و در نبود شناسهٔ بیمار، چشم. دو چشم یک نفر و دو میدان یک چشم هرگز از
دو سوی یک تقسیم قرار نمی‌گیرند.

هیچ‌کدام از سه پیکرهٔ در دسترس شناسهٔ بیمار منتشر نمی‌کنند. پس پیش از نوشتن نخستین تقسیم، برای
هر تصویر یک درهم‌ساز ادراکی و یک توصیفگر الگوی عروقی محاسبه شد، خوشه‌ها ساخته شد، و هر خوشه یک
گروه در نظر گرفته شد. همین کار هم‌زمان تکراری‌های میان‌پیکره‌ای را پیدا کرد، که فرضی نیستند:
\num{۶۸۷} پرونده از \num{۱۷۴۴} پروندهٔ آینهٔ «\lr{Messidor-2}» ما نام‌گذاری به سبک \lr{Messidor-1} دارند.

تقسیم حاصل در \lr{data/splits/} نوشته و کامیت شد و هر مدل این پژوهش آن را از همان‌جا می‌خواند.
اثر انگشت تقسیم \lr{0cfbbfeb081999af} است و در نتایج هر اجرا ثبت می‌شود؛ اجرایی که اثر انگشت
متفاوتی داشته باشد با اجرای دیگر مقایسه نمی‌شود.

\subsection{آنچه کنار گذاشته شد}

دو سطح، چون یک تقسیم پانزده‌درصدی از \lr{IDRiD} چیزی را که دنبالش هستیم اندازه نمی‌گیرد.

\textbf{استخر توسعه} شامل \num{۵۱۶} تصویر \lr{IDRiD} و \num{۱۷۴۴} تصویر \lr{Messidor-2}، جمعاً \num{۲۲۶۰} تصویر
است و آموزش و انتخاب مدل هر دو روی آن انجام می‌شود، با اعتبارسنجی متقاطع گروهی پنج‌تایی.
پیکرهٔ \lr{EyePACS} تنها برای پیش‌آموزش به کار می‌رود و هرگز در انتخاب مدل دخالت نمی‌کند.

\textbf{مجموعهٔ آزمون بیرونی} \lr{APTOS 2019} \cite{aptos2019} است، با \num{۳۶۶۲} تصویر که به‌طور کامل کنار گذاشته
شده و تا انجماد مدل نهایی دیده نمی‌شود. دلیل انتخاب آن این است که از پیش در اختیار بود، پس
کنار گذاشتنش هزینه‌ای نداشت؛ به اندازهٔ کافی بزرگ است که بازهٔ اطمینان حدود یک واحد بدهد؛ و
جابه‌جایی توزیع واقعی است، نه بازبرش همان پیکره.

یک هشدار باید همراه هر عدد بیرونی بیاید: برچسب‌های \lr{APTOS} تک‌داور و نوفه‌دارترند. بخشی از
افت درون‌به‌بیرون نوفهٔ برچسب است، نه ناتوانی در تعمیم، و نباید همه‌اش به مدل نسبت داده شود.

\subsection{انتخاب مدل}

پیکربندی‌ها بر پایهٔ کارایی روی \emph{اعتبارسنجی} رتبه‌بندی می‌شوند، هرگز بر پایهٔ آزمون. این
شامل انتخاب \emph{میان} آزمایش‌ها هم می‌شود، نه فقط فراپارامترها درون یک آزمایش.

این قاعده در نرم‌افزار تثبیت شده است. اسکریپت تولید گزارش، مدل گزارش‌شده را از پروندهٔ
\lr{data/selected\_model.json} می‌خواند و اگر آن پرونده نباشد اجرا را رد می‌کند. پیش‌تر همان
اسکریپت مدل را با بیشینه‌گیری روی \lr{QWK} انتخاب می‌کرد، که انتخاب بر پایهٔ نمره است و وقتی
ارزیابی‌های بیرونی هم در استخر باشند، انتخاب روی مجموعهٔ آزمون می‌شود. اکنون اگر اجرای دیگری
نمرهٔ بالاتری داشته باشد، اسکریپت آن را با صدای بلند اعلام می‌کند و همچنان مدل انتخاب‌شده را
گزارش می‌کند.

\subsection{معیارها}

معیار اصلی برای رتینوپاتی دیابتی، \lr{QWK} است \cite{cohen1968kappa}. این معیار برای خطای دو درجه‌ای بیش از خطای یک
درجه‌ای جریمه می‌گیرد، که همان حالت شکستی است که هدف اصلاح آن را داریم. دقت همیشه در کنارش
گزارش می‌شود و هرگز به تنهایی. بازخوانی هر کلاس نیز همراه هر عدد سرتیتر می‌آید، چون مدل
می‌تواند دقت بگیرد و هم‌زمان نسبت به کلاس نادر کور شود.

برای تصمیم غربالگری، حساسیت و ویژگی رتینوپاتی نیازمند ارجاع گزارش می‌شود. این همان تصمیمی است
که سامانه از آن پشتیبانی می‌کند و همان چیزی است که بالینگر می‌پرسد.

برای ادم ماکولا، تعریف سه‌کلاسهٔ بدون دروازه اصلی است. جزئیات این تصمیم در بخش
\ref{sec:gating} می‌آید.

هر عدد بازهٔ خودگردان‌سازی دارد. چون سطرها همبسته‌اند، \textbf{بازنمونه‌گیری روی گروه‌ها انجام
می‌شود، نه روی تصویرها}. بازنمونه‌گیری روی تصویرها با نزدیک‌تکراری‌ها مانند شواهد مستقل رفتار
می‌کند و بازه را حدود دو برابر باریک‌تر از واقع نشان می‌دهد.

برای مقایسهٔ دو مدل، خودگردان‌سازی جفتی روی همان گروه‌ها انجام می‌شود و بازهٔ اختلاف گزارش
می‌شود. اگر بازه صفر را در بر بگیرد، دو مدل تفکیک‌ناپذیرند و همان‌طور گزارش می‌شوند؛ رتبه‌بندی
نمی‌شوند.

\subsection{واسنجی همتاشده پیش از انتساب}
\label{sec:matched}

\textbf{هیچ اختلافی میان دو مدل به بازنمایی نسبت داده نمی‌شود مگر آن‌که از آزمون واسنجی همتاشده
جان به در ببرد.} پیش از آن‌که بگوییم یک مدل ویژگی‌های بهتری آموخته، هر دو باید به یک شکل واسنجی
شوند — نقاط برش با برازش متقاطع و زیر یک هدف یکسان — و مقایسه تکرار شود.

دلیل وجود این قاعده تجربی است. آستانهٔ پیش‌فرض \num{۰٫۵} روی خروجی سیگموئید چیزی نیست که از مقیاس
درجه‌بندی آمده باشد؛ چیزی است که از تابع فعال‌سازی درمی‌آید. \textbf{یک پیش‌فرض تنظیم‌نشده،
فراپارامتری است که انتخاب شده، نه فراپارامتری که از آن پرهیز شده}، و هر مقایسه‌ای که آن را
میان مدل‌ها ثابت نگه دارد به آن آلوده است.

این قاعده تا امروز چهار ادعا را در \emph{دو جهت} برگردانده و دو بار علامت اختلاف را وارونه
کرده است. سابقه‌اش در جدول \ref{tab:matched_record} فصل بعد آمده است. نکتهٔ مهم این است که
قاعده‌ای که همیشه در یک جهت عمل کند آزمون نیست: سطری که نتیجهٔ مثبت داده همان چیزی است که به
سطرهای دیگر اعتبار می‌دهد.

\subsection{دو صلاحیت که همراه هر ادعای کلاسی می‌آیند}

نخست، \textbf{هیچ تصحیحی برای مقایسه‌های چندگانه اعمال نمی‌شود}. پنج کلاس ضرب در شرایط
مقایسه‌شده یعنی به‌طور متوسط یک خانهٔ معنادار از سر تصادف انتظار می‌رود. پس اعتبار از
\emph{تکرار یک اختلاف در شرایط مختلف} می‌آید، نه از مقدار \lr{p} به تنهایی.

دوم، جایی که مقایسه بیش از یک چیز را عوض می‌کند، ادعای کلاسی همهٔ مخدوش‌کننده‌های ادعای کلی را
به ارث می‌برد.

\section{داده}

\subsection{پیکره‌ها}

\lr{IDRiD} \cite{porwal2018idrid} با \num{۵۱۶} تصویر تنها منبع سه‌کلاسهٔ ادم ماکولا در این پژوهش است و
\textbf{\num{۵۱} تصویر در درجهٔ میانی} دارد. همین عدد سقف کارایی سر ادم ماکولا را تعیین می‌کند و در
فصل \ref{ch:results} به آن بازمی‌گردیم.

\lr{Messidor-2} \cite{decenciere2014messidor} با \num{۱۷۴۴} تصویر برچسب رتینوپاتی دارد و برای ادم ماکولا تنها برچسب دودویی
«نیازمند ارجاع یا نه» می‌دهد. این برچسب جزئی است، نه ناقص: در صورت‌بندی آستانه‌ای، «نیازمند
ارجاع» دقیقاً معادل درجهٔ دوی \lr{IDRiD} است، پس آستانهٔ دوم را دقیق نظارت می‌کند و آستانهٔ
نخست را پوشانده می‌گذارد. هیچ تقریبی در کار نیست.

\lr{APTOS 2019} \cite{aptos2019} با \num{۳۶۶۲} تصویر به‌طور کامل کنار گذاشته شده است.

\subsection{ارزیابی بدون دروازه به‌عنوان تعریف اصلی}
\label{sec:gating}

کار پیشین شاخهٔ ادم ماکولا را مشروط به رتینوپاتی بزرگ‌تر یا مساوی یک می‌کرد. در این پژوهش
ارزیابی \textbf{بدون دروازه} اصلی است و عدد مشروط به‌عنوان ثانویه و همیشه برچسب‌خورده گزارش
می‌شود. سه دلیل دارد.

نخست، \textbf{دروازه چیزی را کنار نمی‌گذارد}. روی \lr{IDRiD} بررسی شد: از \num{۵۱۶} تصویر، هر \num{۱۶۸}
تصویری که رتینوپاتی صفر دارند، درجهٔ ادم ماکولایشان هم صفر است. هیچ نمونهٔ مثبت ادم ماکولا با
رتینوپاتی صفر وجود ندارد. پس دروازه از چیزی محافظت نمی‌کند و تنها مخرج کسر را عوض می‌کند.

دوم، معیار مشروط به درست بودن سر رتینوپاتی وابسته است، پس میان مدل‌هایی با سرهای متفاوت
قابل مقایسه نیست.

سوم، و مهم‌تر از همه، \textbf{کف کلاس اکثریت به‌شدت فرق می‌کند}: \num{۶۹٫۸} درصد برای مجموعهٔ مشروط
در برابر \num{۴۷٫۱} درصد برای بدون دروازه. \textbf{هر عدد ادم ماکولا در این پژوهش کنار کف خودش نقل
می‌شود.}

\section{معماری و آموزش}

ستون فقرات مشترک است و دو سر جدا دارد، یکی برای رتینوپاتی با پنج درجه و یکی برای ادم ماکولا با
سه درجه. سرها آستانه‌ای‌اند: به جای بیشینهٔ نرم روی \lr{K} کلاس، \lr{K-1} احتمال از شکل
$P(y>k)$ تخمین زده می‌شود. این انتخاب سه کار می‌کند: ترتیب درجه‌ها را وارد تابع زیان می‌کند
که همان چیزی است که \lr{QWK} می‌سنجد؛ رمزگشایی را با ساخت سازگار با رتبه می‌کند؛ و از همه
مهم‌تر، برچسب جزئی \lr{Messidor-2} را قابل استفاده می‌کند.

درجه‌بندی پنج‌درجه‌ای رتینوپاتی و سه‌درجه‌ای ادم ماکولا از مقیاس بالینی بین‌المللی
می‌آید \cite{wilkinson2003scale}. ستون فقرات پیش‌فرض \lr{DenseNet121}
\cite{huang2017densenet} است و مدل انتخاب‌شده \lr{EfficientNet-B3}
\cite{tan2019efficientnet}؛ \lr{ConvNeXt-tiny} \cite{liu2022convnext} نیز به‌عنوان یک نامزد
آزموده شد. دو سر ترتیبی جایگزین که در فصل \ref{ch:results} ابطال می‌شوند
\lr{CORAL} \cite{cao2020rank} و \lr{CORN} \cite{shi2023corn} هستند، و راهبرد
\lr{LP-FT} از \cite{kumar2022lpft} می‌آید. گزارش این پژوهش در برابر سیاههٔ
\lr{CLAIM} \cite{mongan2020claim} و \lr{TRIPOD} \cite{collins2015tripod} سنجیده شده است، و
معیار واسنجی و مسئلهٔ بیش‌اطمینانی شبکه‌های عمیق از \cite{guo2017calibration} گرفته شده است.

پیش‌آموزش روی \lr{EyePACS} \cite{eyepacs2015} انجام می‌شود، یک بار، و همهٔ تاها از همان وزن
آغاز می‌کنند.

افزایش داده شامل چرخش آزاد، بازتاب، و لرزش ملایم رنگ و گاما است. چرخش نامحدود است چون شبکیه
جهت متعارف ندارد و هر معیار درجه‌بندی نسبت به چرخش ناوردا است.

\section{تبار داده}
\label{sec:provenance}

این بخش سهمی است، نه اعتراف. شش اشکال تبار داده در این پژوهش هیچ خطای زمان اجرا تولید نکردند و
در سکوت عدد غلط ساختند. سازوکارهای زیر پاسخ آن‌هاست.

\subsection{پیکربندی، مصرف نیست}

چهار سازوکار مستقل تبار داده هم‌زمان از اجرایی عبور کردند که ترکیبش غلط می‌بود: بلوک پیکربندی
یکسان، اثر انگشت تقسیم یکسان، نگهبان اسکریپت، و میدانی که برای هر دو اجرا \emph{به‌راستی درست}
بود و برای هرکدام معنای دیگری داشت. هر چهار سازوکار ثبت می‌کنند که اجرا چه چیزی را
\emph{پیکربندی} کرده است؛ هیچ‌کدام ثبت نمی‌کنند چه چیزی را \emph{مصرف} کرده است.

پس هر اجرا اکنون یک \textbf{بیانیهٔ مصرف} می‌نویسد: برای هر پیکره، تعداد تصویرهای بازیابی‌شده و
یک درهم‌ساز روی مسیرهای بازیابی‌شده نسبت به محل سوارشدن دیتاست. دو اجرا تنها وقتی ترکیب،
هم‌گروه یا مقایسه می‌شوند که بیانیه‌هایشان بخواند. مقایسه‌های عمداً میان‌منبعی ممکن‌اند، اما
پرچم لازم اختلاف را \emph{در سند خروجی مهر می‌کند}؛ هشداری که خواننده نبیند هشدار نیست.

\subsection{پنج نمونه، که یکی از آن‌ها خودِ ماست}
\label{sec:provenance-exhibits}

استدلال این بخش وقتی وزن دارد که نمونه داشته باشد. چهار نمونه از کارهای منتشرشده می‌آید و
\textbf{نمونهٔ پنجم از همین پایان‌نامه است}. ترتیب عمدی است: فصلی که تبار دادهٔ دیگران را
حسابرسی کند و شکست خودش را از همان جنس گزارش نکند، ادعایی می‌کند که خودش را زیر آن نمی‌برد.

\textbf{نمونهٔ یکم — سطح بالای \lr{AUC} به‌جای واسنجی \cite{yu2026dualswinord}.} یک مقالهٔ
داوری‌شدهٔ \num{۲۰۲۶} می‌نویسد که مقادیر بالای \lr{AUC} «تأیید می‌کنند» مدل احتمال‌های
خوش‌واسنجی می‌دهد. آن مقاله \emph{هیچ} معیار واسنجی گزارش نمی‌کند: نه \lr{ECE}، نه
\lr{Brier}، نه منحنی اطمینان. و \lr{AUC} نسبت به هر تبدیل یکنوای نمره‌ها ناوردا است، پس
دربارهٔ واسنجی هیچ نمی‌گوید؛ مدلی می‌تواند همان \lr{AUC} را نگه دارد و به هر اندازه
بدواسنجی باشد.

\textbf{نمونهٔ دوم — تحلیل چشمی روی داده‌های جفت‌شده \cite{tomic2024handheld}.} پژوهشی با \num{۳۲۰}
چشم از \num{۱۶۰} بیمار، همهٔ تحلیل‌ها را چشم‌به‌چشم انجام می‌دهد و برای همبستگی دو چشم یک نفر تصحیحی
نمی‌گذارد. بازه‌های اطمینان گزارش‌شده از این راه باریک‌تر از واقع درمی‌آیند. این همان دلیلی است
که در بخش \ref{sec:matched} بازنمونه‌گیری این پژوهش روی گروه‌ها انجام می‌شود، نه روی تصویرها.

\textbf{نمونهٔ سوم — ستونی که از ستون دیگر رونویسی شده \cite{tahir2025meta}.} در یک
فراتحلیل، ستون منفی‌های کاذب عیناً تکرار ستون مثبت‌های کاذب است. این خطا در جدول نمی‌ماند: از
همان اعداد، حساسیت یکی از پژوهش‌های واردشده \num{۰٫۰۶} درمی‌آید، در حالی که حساسیت تجمیعی گزارش‌شده
\num{۰٫۹۵} است. \textbf{خرابی داده به عددهای گزارش‌شده نشت کرده است.}

\textbf{نمونهٔ چهارم — تقسیمی که در دو پیمانه اعمال می‌شود و در سومی نه
\cite{tang2026mmrdr}.} یک توصیف‌نامهٔ داده در خانوادهٔ \lr{Nature} زیرمجموعه‌های خود را در دو
پیمانه در سطح بیمار تقسیم می‌کند «برای جلوگیری از نشت داده»، و زیرمجموعهٔ سوم را در سطح تصویر،
چون پیکرهٔ مبدأ شناسهٔ بیمار ندارد؛ سپس یک جدول محک واحد روی هر سه گزارش می‌شود. قاعده
شناخته شده، نوشته شده، و در یک‌سوم موارد اجرا نشده است.

\subsection{نمونهٔ پنجم: شکست خودِ این متن}
\label{sec:our-own-failure}

فصل حاضر در نسخهٔ نخست خود با این جمله آغاز می‌شد که همهٔ عددهایش از اجراهای آرشیوشده می‌آیند و
هیچ عددی با دست تایپ نشده است. \textbf{این جمله درست نبود.} عددها با دست تایپ شده بودند، هیچ
سازوکاری آن‌ها را بررسی نمی‌کرد، و سه عدد غلط بود.

\begin{table}[H]
\caption{سه خطای واقعی در نسخه‌های نخست همین متن، که تنها به این دلیل پیدا شدند که ساختن یک
دفتر اعداد، حل کردن هر رقم در برابر یک مصنوع آرشیوی را ضروری کرد.}
\centering
\begin{tabular}{|l|c|c|}
\hline
\textbf{عدد} & \textbf{نوشته‌شده} & \textbf{درست} \\ \hline
کف مشروط ادم ماکولا & \num{۶۹٫۶}٪ & \num{۶۹٫۸}٪ \\ \hline
تطابق دقیق در فصل \ref{sec:parta} & \num{۹۶٫۶}٪ & \num{۹۶٫۳}٪ \\ \hline
اختلاف هم‌نهاد تصویرها & \num{۰٫۷۸} تا \num{۰٫۹۶} & \num{۰٫۶۵} تا \num{۰٫۷۵} \\ \hline
\end{tabular}
\label{tab:our-failure}
\end{table}

\textbf{سطر سوم یافته است، نه غلط تایپی.} آن بازه شاهدِ نخستین توضیح رقیبی بود که فصل
\ref{sec:parta} باید رد می‌کرد: این‌که تناظر میان شماره‌گذاری قطعه‌بندی و شماره‌گذاری
درجه‌بندی، ماسک را کنار درجهٔ اشتباه گذاشته باشد. تمام آن حکم روی این عدد می‌ایستاد.

\textbf{و آن عدد هیچ مصنوعی نداشت.} نه اسکریپتی، نه پروندهٔ خروجی، نه ثبتی از روش. بازتولید
نشد، و وقتی از نو محاسبه شد پاسخ \num{۰٫۶۵} تا \num{۰٫۷۵} درآمد. نتیجه عوض نشد — جفت‌ها همان
عکس‌اند و حکم سرجایش است — ولی \textbf{عددهای منتشرشده بازتولیدپذیر نبودند و با عددهایی
جایگزین شدند که هستند.} عددی که اسکریپتش دیگر وجود ندارد، ادعا است نه اندازه‌گیری؛ و این همان
جمله‌ای است که این فصل دربارهٔ مقاله‌های دیگران می‌گوید.

\subsection{و نمونهٔ ششم، که هنگام نوشتن نمونهٔ پنجم ساخته شد}

سه خطای بالا هنگام بررسی گسترده‌تری پیدا شدند، و آن بررسی خودش به یک اشتباه انجامید که آموزندهٔ
هر پنج نمونهٔ پیشین است.

مشاهده این بود که عددهای اعشاری فصل‌های نخست، بخش صحیح و بخش اعشاری‌شان جابه‌جا نوشته شده است.
تشخیص از روی استاندارد ساخته شد: ارقام فارسی در جدول یونیکد از ردهٔ عدد اروپایی‌اند، رشته‌های
عددی چپ‌به‌راست چیده می‌شوند، پس متن همان‌طور که نوشته شده چاپ می‌شود، پس
\textbf{فصل روش به خواننده می‌گفت آستانهٔ پیش‌فرض \wrongnum{۵/۰} است}.

\textbf{استدلال دربارهٔ یونیکد درست بود و نتیجه دربارهٔ این سند غلط.} پیاده‌سازی دوسویهٔ
موتور حروف‌چینی این پایان‌نامه، پیاده‌سازی کامل الگوریتم یونیکد نیست. عددی که به ترتیب منطقی
نوشته شود، \emph{وارونه} چاپ می‌شود. \textbf{پس نسخه‌های نخست درست چاپ می‌شدند}؛ نویسنده‌شان
برای موتور چاپ می‌نوشت. «تصحیح» منبع را منطقی کرد و \textbf{صفحهٔ چاپ‌شده را غلط}، و همین‌طور
ماند تا وقتی که سرانجام یک صفحه چاپ و نگاه شد.

\textbf{جنس این خطا دقیقاً همان جنس چهار نمونهٔ منتشرشده است.} در نمونهٔ نخست، ادعای واسنجی
جای اندازه‌گیری واسنجی را گرفته بود. اینجا، مشخصات یک استاندارد جای نگاه کردن به خروجی را
گرفت. \textbf{هیچ‌کدام به خودِ چیز نگاه نکردند.}

قرارداد عوض شد، ولی \emph{به دلیل خودش}، نه به‌عنوان رفع اشکال: منبع به ترتیب منطقی نوشته
می‌شود و هر عدد در یک بازهٔ صریح چپ‌به‌راست پیچیده می‌شود. این هم درست چاپ می‌شود و هم
جست‌وجوپذیر و بررسی‌پذیر است؛ قرارداد پیشین درست چاپ می‌شد و هیچ‌کدام از این دو نبود.

\subsection{آنچه اکنون این را نگه می‌دارد}

هر عدد فارسیِ متن باید در \lr{docs/generated/thesis\_numbers.json} باشد، که از مصنوع‌های
آرشیوی ساخته می‌شود. ویراستار خودکار سه چیز را رد می‌کند: رقمی که در آن دفتر نباشد، جداکنندهٔ
اعشار نادرست، و \textbf{عددی که در پوشش چپ‌به‌راست نباشد}. آزمون سوم دقیقاً به این دلیل وجود
دارد که نمونهٔ ششم از دو آزمون نخست عبور کرد.

هر سه آزمون پیش از اعتماد، در هر دو جهت عمل کرده‌اند. \textbf{آنچه ثابت نمی‌کنند} این است که
عدد برای \emph{آن جمله} درست باشد؛ هیچ آزمون خودکاری آن را ثابت نمی‌کند، و ادعای خلافش دقیقاً
همان «آزمونی که موفقیت گزارش می‌کند بی‌آن‌که چیزی را بررسی کرده باشد» است.

\textbf{درسی که می‌ماند ساده است:} ادعایی دربارهٔ خروجی، با نگاه کردن به خروجی حل می‌شود. دو
نگهبان این پژوهش به این دلیل وجود دارند که یک بار استدلال جای بازرسی را گرفت.

\subsection{آزمونی که فقط در یک جهت عمل کند آزمون نیست}

هر قاعده در این پژوهش سابقه‌اش را حمل می‌کند، و سابقه باید ابطال‌هایی در هر دو جهت داشته باشد.
قاعده‌ای که تنها نتایج ناخوشایند را حل کرده باشد از «راهی برای توجیه» قابل تشخیص نیست، و
قاعده‌ای که تنها تأیید کرده باشد مهر لاستیکی است.

حالت انحطاطی این وضع بدترین شکست ممکن برای یک نگهبان است: \textbf{آزمونی که موفقیت گزارش کند
بی‌آن‌که چیزی را بررسی کرده باشد.} دو نمونه از آن در این پژوهش رخ داد و هر دو اکنون ادعا
می‌کنند که چیزی بازرسی شده است. نگهبانی که با صدای بلند شکست بخورد یک اجرا هزینه دارد؛
نگهبانی که به غلط قبول شود، اعتماد به این‌که اصلاً بررسی کرده‌اید را هزینه دارد.
